\documentclass[12pt]{article}
\usepackage[T1]{fontenc}
\usepackage[utf8]{inputenc}
\usepackage{lmodern}
\usepackage{textcomp}
\usepackage{enumitem}
\usepackage{geometry}
\geometry{margin=1in}
\begin{document}
\begin{center}
Answer Key - Chemistry - Undergraduate Level Assessment 23 \
\small (Answers are ordered exactly as in the question paper.)
\end{center}

\section*{Multiple Choice}
\begin{enumerate}[label=\arabic*.]
\item \textbf{Answer:} B
\\ \textit{Options:} A) Ti$_{3}$+; B) Cr$_{3}$+; C) Fe$_{3}$+; D) Zn$_{2}$+
\\ \textit{Note:} Cr$_{3}$+ is a d 3 transition metal ion that has no unpaired electrons in its ground state, making it diamagnetic and unsuitable as a paramagnetic quencher.
\item \textbf{Answer:} C
\\ \textit{Options:} A) 3; B) 2; C) 1; D) 0
\\ \textit{Note:} The orbital angular momentum quantum number, l, of 1 corresponds to the p-orbital electrons in aluminum, which are higher in energy and more loosely held compared to the s-orbital electrons, making them the most easily removed during ionization.
\item \textbf{Answer:} C
\\ \textit{Options:} A) K+; B) Ca$_{2}$+; C) Sc$_{3}$+; D) Rb+
\\ \textit{Note:} Sc$_{3}$+ has the smallest radius because it has lost three electrons, resulting in a higher effective nuclear charge relative to its remaining electrons, which pulls them closer to the nucleus and decreases the ionic radius.
\item \textbf{Answer:} C
\\ \textit{Options:} A) Nuclear magnetic resonance; B) Infrared; C) Raman; D) Ultraviolet-visible
\\ \textit{Note:} Raman spectroscopy is a light-scattering technique because it relies on the inelastic scattering of monochromatic light, typically from a laser, to provide information about molecular vibrations and chemical composition.
\item \textbf{Answer:} C
\\ \textit{Options:} A) all molecular bonds absorb IR radiation; B) IR peak intensities are related to molecular mass; C) most organic functional groups absorb in a characteristic region of the IR spectrum; D) each element absorbs at a characteristic wavelength
\\ \textit{Note:} The correct answer is based on the fact that different organic functional groups exhibit unique absorption patterns in specific regions of the infrared spectrum, allowing chemists to identify and characterize them based on their distinct IR signatures.
\end{enumerate}

\section*{True / False}
\begin{enumerate}[label=\arabic*.]
\setcounter{enumi}{5}
\item \textbf{Answer:} True
\\ \textit{Options:} A) True; B) False
\\ \textit{Note:} The atomic radii of the lanthanide elements decrease from La to Lu because the increasing nuclear charge from the additional protons attracts the electrons more strongly, resulting in a greater effective nuclear charge that pulls the electron cloud closer to the nucleus.
\item \textbf{Answer:} True
\\ \textit{Options:} A) True; B) False
\\ \textit{Note:} This statement is true because of the Heisenberg Uncertainty Principle, which asserts that the more precisely one property (such as position) is measured, the less precisely the other property (momentum) can be known.
\item \textbf{Answer:} False
\\ \textit{Options:} A) True; B) False
\\ \textit{Note:} A magnetic field of 0.5 T is insufficient to significantly polarize 13 C at 298 K due to its low nuclear magnetic moment and the thermal energy at that temperature, which limits the degree of polarization achievable.
\item \textbf{Answer:} True
\\ \textit{Options:} A) True; B) False
\\ \textit{Note:} The statement is true because, according to the Gibbs phase rule, the number of phases (P) in equilibrium in a system can be determined using the equation P = C - F + 2, where C is the number of components and F is the number of degrees of freedom, and for a three-component system at a fixed temperature and pressure, a maximum of five phases can coexist.
\item \textbf{Answer:} False
\\ \textit{Options:} A) True; B) False
\\ \textit{Note:} The statement is false because, although NaCl is an ionic compound, MgO has a greater lattice enthalpy due to the higher charges of the magnesium and oxide ions, which results in stronger electrostatic attractions compared to the sodium and chloride ions.
\end{enumerate}

\section*{Long Form Response}
\begin{enumerate}[label=\arabic*.]
\setcounter{enumi}{10}
\item \textbf{Answer:} The R-branch in the vibrational-rotational spectrum of a diatomic molecule is significant because it corresponds to transitions where the rotational quantum number increases by one (ΔJ = 1) while the vibrational quantum number also increases by one (Δu = 1). These transitions imply that as the molecule absorbs energy, it not only moves to a higher vibrational state but also gains rotational energy, reflecting the coupling between vibrational and rotational motions. The R-branch can provide insights into the moments of inertia and bond lengths of the molecule, allowing chemists to infer structural information and understand molecular behavior under vibrational excitation. Overall, the analysis of the R-branch offers essential information about the energy levels and dynamics of diatomic molecules.
\\ \textit{Note:} The answer correctly identifies that the R-branch represents transitions in which both vibrational and rotational quantum numbers increase, illustrating the coupling of vibrational and rotational motions and their significance in understanding molecular energy levels and behavior.
\item \textbf{Answer:} The NMR frequency of 31 P in a 20.0 T magnetic field can be calculated using the formula ν = γB 0/(2π), where ν is the frequency, γ is the gyromagnetic ratio for 31 P (approximately 108.3 MHz/T), and B$_{0}$ is the magnetic field strength in teslas. Substituting the values, we find ν = 108.3 MHz/T \ensuremath{\times} 20.0 T / (2π) \ensuremath{\approx} 345.0 MHz. The frequency of NMR signals is influenced by factors such as the strength of the magnetic field, the gyromagnetic ratio of the nucleus being analyzed, and the chemical environment surrounding the nuclei, which can lead to shifts in resonance frequency due to electron shielding or deshielding effects. Understanding these factors is crucial for interpreting NMR spectra in chemical analysis.
\\ \textit{Note:} correct because it accurately applies the formula for calculating the NMR frequency by incorporating the gyromagnetic ratio of 31 P and the strength of the magnetic field, while also acknowledging the key factors that influence NMR frequencies, such as magnetic field strength and chemical environment.
\end{enumerate}

\vfill
\noindent\textit{End of Answer Key}
\end{document}